\begin{Theorem}
\label{thm:sobolev-banach-separability}
The space $W^{k,p}(\Omega)$ is a Banach space. If $p<\infty$, it is separable.
\end{Theorem}

\begingroup
\renewcommand{\endProof}{\par\nobreak\hfill$\blacksquare$\par\vspace{0.1cm}}
\begin{Proof}
For $u\in W^{k,p}(\Omega)$, let $u_\alpha=D^\alpha u$. Then the mapping $u\to(u_\alpha)_{|\alpha|\leq k}$ is an isometry from $W^{k,p}(\Omega)$ onto a subspace $M$ of the direct product $\prod_{|\alpha|\leq k}L^p(\Omega)$. The completeness of $W^{k,p}(\Omega)$ follows if we show that $M$ is closed. For this, one has to show that if $u_n\in W^{k,p}$ and $D^\alpha u_n\to v^\alpha$ in $L^p(\Omega)$, then $v^\alpha=D^\alpha v^0$. But this is clear, since $L^p$-convergence implies distributional convergence.

If $p<\infty$, then $\prod_{|\alpha|\leq k}L^p(\Omega)$ is separable. Hence we only need to show that a subspace of a separable Banach space is separable. This is done in the following lemma to complete the proof of this theorem.
\end{Proof}
\endgroup

\begin{Lemma}
\label{lem:separable-banach-linear-subspace}
Let $X$ be a separable Banach space and let $M$ be a linear subspace of $X$. Then $M$ is separable.
\end{Lemma}

\begin{Proof}
For any given $\varepsilon>0$, there is a countable set $\{x_n\}_{n\in\mathbb{N}}$ in $X$ such that for any $x\in X$ there is $n\in\mathbb{N}$ with $\|x-x_n\|\leq\varepsilon$. From the sequence $x_n$, discard all elements such that the ball of radius $\varepsilon$ centered at $x_n$ does not intersect $M$. The remaining points form a new sequence $y_n$. Now let $u_n$ be any point in $M$ which lies in the $\varepsilon$-neighborhood of $y_n$. Then every point in $M$ lies within $\varepsilon$ of one of the $y_n$ and hence within $2\varepsilon$ of one of the $u_n$. By letting $\varepsilon$ run through any sequence that converges to zero, we can generate a countable dense set in $M$.
\end{Proof}

If $p<\infty$, then $\mathcal{D}(\Omega)$ is dense in $L^p(\Omega)$. This is generally not the case for the spaces $W^{k,p}(\Omega)$. Indeed we have the following result.

\begin{Lemma}
\label{lem:test-functions-sobolev-nondensity}
Let $\Omega$ be bounded, let $k\geq1$ and $1\leq p\leq\infty$. Then $\mathcal{D}(\Omega)$ is not dense in $W^{k,p}(\Omega)$.
\end{Lemma}

\begin{Proof}
If $u$ is any function in $C^{\infty}(\overline\Omega)$, then the mapping
\begin{equation}\label{eq:sobolev-test-functions-annihilating-functional}
L:\phi\mapsto\int_{\Omega}\sum_{|\alpha|\leq k}D^\alpha\bar u(\bm{x})D^\alpha\phi(\bm{x})\,\mathrm{d}\bm{x}
\end{equation}
defines a bounded linear functional on $W^{k,p}(\Omega)$. Moreover, if
\begin{equation}\label{eq:sobolev-nondensity-annihilating-elliptic-equation}
\sum_{|\alpha|\leq k}(-1)^{|\alpha|}D^{2\alpha}\bar u=0,
\end{equation}
then $L(\phi)$ vanishes for every $\phi\in\mathcal{D}(\Omega)$. A nonzero solution $u$ of \eqref{eq:sobolev-nondensity-annihilating-elliptic-equation} in any parallelepiped containing $\Omega$ can easily be found by separation of variables; see Chapter~\ref{chap:introduction}. Moreover, $L(u)>0$ and $u\in W^{k,p}(\Omega)$; hence $\mathcal{D}(\Omega)$ cannot be dense.
\end{Proof}

The last lemma motivates the following definition.

\begin{Definition}
\label{def:sobolev-test-function-closure}
By $W_0^{k,p}(\Omega)$ we denote the closure of $\mathcal{D}(\Omega)$ in $W^{k,p}(\Omega)$.
\end{Definition}
